\chapter*{第 5 章 管路、孔口、管嘴的水力计算}\addcontentsline{toc}{chapter}{第 5 章 管路、孔口、管嘴的水力计算}\markboth{第 5 章 管路、孔口、管嘴的水力计算}{}

\begin{tishi}
\textbf{本章是全书篇幅最长、计算题命中率最高的一章}（E1 slide33「第五章 孔口/管嘴/管路水力计算＝考试计算题重点」；
E2「计算题占 50\% 分值，第五章同为主要计算考点」；E3「每年的真题大部分都源于课后习题的改编」）。
教材原书本章习题共 \textbf{20 题}（5.1--5.20，书 p95--98），本册\textbf{一题不漏}全部收录。\\
做题要求：计算题必须写出\textbf{已知 $\to$ 求 $\to$ 公式 $\to$ 代入 $\to$ 结果 $\to$ 单位检验}六段；
\textbf{凡涉及流态或分区的题，第一步必须先算 $\Rey$ 并写清判据}（这是本章最大的过程分来源）。
本章的判据与系数属于"背成条件反射"的内容：$\Rey<2000$ 层流、$\lambda=64/\Rey$；
孔口 $\mu\approx0.62$、圆柱形外管嘴 $\mu\approx0.82$；进口 $\zeta=0.5$、出口 $\zeta=1.0$；
管嘴正常条件 $H_0\le9\ \mathrm{m}$ 且 $l=(3\sim4)d$；串联"流量相等、损失相加"、并联"损失相等、流量相加"。\\
教材原书习题图一律用〔图示〕文字精确转述，\textbf{不重绘图、不编造尺寸}；OCR 明显错字已按量纲订正，逐题列在每题后的〔订正〕条中。
\end{tishi}

\begin{ti}
\sloppy
\emergencystretch=2em

\item 选择题。\xstarfull{5}\yiju{E4 2015 年试题选择题第 10 题（原文即"圆管层流流动，过流断面上切应力分布为"，与本题第（1）小问一字不差）；E4 2015 年试题选择第 11 题（"雷诺数的物理意义"＝本题第（3）（4）小问的判据来源）；E4 2015 年试题填空第 5、6 题（层流断面平均速度、$\Rey$ 判流态）；E1 slide33「第五章＝考试计算题重点」；E5 题库 \texttt{10.pdf} 选择第 10 题（圆管与方管沿程损失比 0.886）}\nandu{中}

\begin{choices}
\item 圆管层流流动，过流断面上切应力分布为（\hspace{1cm}）。
\item[] A. 在过流断面上是常数\quad B. 管轴处是零，且与半径成正比\quad C. 管壁处是零，向管轴线性增大\quad D. 按抛物线分布
\item 在圆管流中，层流的断面流速分布符合（\hspace{1cm}）。
\item[] A. 均匀规律\quad B. 直线变化规律\quad C. 抛物线规律\quad D. 对数曲线规律
\item 圆管紊流过渡区的沿程阻力系数（\hspace{1cm}）。
\item[] A. 与雷诺数 $\Rey$ 有关\quad B. 与管壁相对粗糙度 $\Delta/d$ 有关\quad C. 与 $\Rey$ 和 $\Delta/d$ 有关\quad D. 与 $\Rey$ 和管长 $L$ 有关
\item 圆管紊流粗糙区的沿程阻力系数（\hspace{1cm}）。
\item[] A. 与雷诺数 $\Rey$ 有关\quad B. 与管壁相对粗糙度 $\Delta/d$ 有关\quad C. 与 $\Rey$ 和 $\Delta/d$ 有关\quad D. 与 $\Rey$ 和管长 $L$ 有关
\item 水流在管道直径、水温、沿程阻力系数都一定时，随着流量的增加，黏性底层的厚度就（\hspace{1cm}）。
\item[] A. 增加\quad B. 减小\quad C. 不变\quad D. 不定
\item 串联管道各串联管段的（\hspace{1cm}）。
\item[] A. 水头损失相等\quad B. 总能量损失相等\quad C. 水力坡度相等\quad D. 所通过的流量相等
\item 长管并联管道各并联管段的（\hspace{1cm}）。
\item[] A. 水头损失相等\quad B. 总能量损失相等\quad C. 水力坡度相等\quad D. 通过的水量相等
\item 并联长管 1、2，两管的直径、沿程阻力系数均相同，长度 $L_2=3L_1$，则通过的流量为（\hspace{1cm}）。
\item[] A. $Q_1=Q_2$\quad B. $Q_1=1.5Q_2$\quad C. $Q_1=1.73Q_2$\quad D. $Q_1=3Q_2$
\item 两水池水位差为 $H$，用两根等径等长、沿程阻力系数均相同的管道连接，如习题 5.1（9）图所示，按长管考虑，则（\hspace{1cm}）。
\item[] A. $Q_1>Q_2$\quad B. $Q_1=Q_2$\quad C. $Q_1<Q_2$\quad D. $Q_2=0$
\item 在正常工作条件下，作用水头 $H$、直径 $d$ 相等时，小孔口的流量 $Q$ 与圆柱形直角外管嘴的流量 $Q_n$ 相比较，有（\hspace{1cm}）。
\item[] A. $Q>Q_n$\quad B. $Q=Q_n$\quad C. $Q<Q_n$\quad D. 不定
\item 孔口出流的流量系数、流速系数、收缩系数从小到大的正常排序是（\hspace{1cm}）。
\item[] A. 流量系数、流速系数、收缩系数\quad B. 流量系数、收缩系数、流速系数\quad C. 流速系数、流量系数、收缩系数\quad D. 收缩系数、流量系数、流速系数
\end{choices}

\item 两种液体阻力及能量损失形式和它们的计算公式分别是什么？\xstarfull{5}\yiju{E1 slide16「简答考水头损失分类及定义」；E1 slide33「第五章＝考试计算题重点」；E2「第五章同为主要计算考点」；E4 2026 回忆版简答考"水击现象、层流与雷诺数关系"（同属本章概念题）}\nandu{易}

\item 一等径圆管内径 $d=100\ \mathrm{mm}$，流通运动黏度 $\nu=1.306\times10^{-6}\ \mathrm{m^2/s}$ 的水，求管中保持层流流态的最大流量 $Q$。\xstarfull{4}\yiju{E4 2022 年 818 单选第 9 题（$d=10\ \mathrm{mm}$、$30\ \mathrm{mm}$ 两圆管判流态，同属"$\Rey$ 判流态"考点，方向相反）；E4 2015 年试题填空第 6 题（$\Rey$ 判流态）；E1 slide33「第五章＝考试计算题重点」}\nandu{中}

\item 证明圆管层流过流断面的流速分布为 $u=\dfrac{\Delta p}{4\mu L}\left(R^2-r^2\right)$，其中 $L$ 为管长，$R$ 为管道半径，$\Delta p$ 为压差，$\mu$ 为动力黏度。\xstarfull{4}\yiju{E4 2015 年试题选择题第 10 题（"圆管层流流动，过流断面上切应力分布为"，与本题的速度抛物线分布互为表里，同属教材 5.3 节）；E4 2015 年试题填空第 5 题（"圆管的流动为层流时，断面平均速度是最大速度的\ \rule{1.2cm}{0.4pt}\ 倍"，正是本题结论的直接推论）；E1 slide33「第五章＝考试计算题重点」}\nandu{中}

\item 利用毛细管测定油液黏度，已知毛细管直径 $d=4.0\ \mathrm{mm}$，长度 $L=0.5\ \mathrm{m}$，流量 $Q=1.0\ \mathrm{cm^3/s}$，测压管落差 $h=15\ \mathrm{cm}$，如习题 5.5 图所示。管中作层流流动，求油液的运动黏度 $\nu$。
\tuyi{习题 5.5 图：一段水平放置的\textbf{毛细管}（细玻璃管），管内油液作层流流动；毛细管\textbf{直径 $d=4.0\ \mathrm{mm}$}、\textbf{长度 $L=0.5\ \mathrm{m}$}（图中以 $d$ 标注管径、以 $L$ 标注测量段长度）；毛细管两端就近各接出一根竖直的\textbf{测压管}（图中以两根竖直细管表示），两测压管内的油液面\textbf{高度差为 $h=15\ \mathrm{cm}$}（图中以 $h$ 标注在两液面之间）；通过毛细管的流量为 $Q=1.0\ \mathrm{cm^3/s}$；测压管落差 $h$ 对应的压差全部用于克服管段的沿程阻力，局部阻力不计。}
\xstarfull{4}\yiju{E4 2015 年试题选择题第 10 题与填空第 5 题（同属"圆管层流"考点）；E1 slide33「第五章＝考试计算题重点」；E6 教材 5.3 节哈根--泊肃叶流量公式 $Q=\pi d^4\Delta p/(128\mu L)$ 的直接应用}\nandu{中}

\item $\rho=0.85\ \mathrm{g/cm^3}$、$\nu=0.18\ \mathrm{cm^2/s}$ 的油在管径为 $100\ \mathrm{mm}$ 的管中以 $u=6.35\ \mathrm{cm/s}$ 的速度作层流运动，求：（1）管中心处的最大流速；（2）在离管中心 $r=20\ \mathrm{mm}$ 处的流速；（3）沿程阻力系数 $\lambda$；（4）管壁切应力 $\tau_0$ 及每 $1000\ \mathrm{km}$ 管长的水头损失。\xstarfull{4}\yiju{E4 2015 年试题选择题第 10 题（层流切应力分布：管轴处为零、与半径成正比）；E4 2015 年试题填空第 5 题（层流断面平均速度＝最大速度的 0.5 倍）；E1 slide33「第五章＝考试计算题重点」}\nandu{中}

\item 应用细管式黏度计测定油的黏度，已知细管直径 $d=6\ \mathrm{mm}$，测量段长 $L=2\ \mathrm{m}$，如习题 5.7 图所示。实测油的流量 $Q=77\ \mathrm{cm^3/s}$，水银压差计的读值 $h=30\ \mathrm{cm}$，油的密度 $\rho=900\ \mathrm{kg/m^3}$。试求油的运动黏度 $\nu$。
\tuyi{习题 5.7 图：一段水平放置的\textbf{细管}（黏度计测量段），\textbf{直径 $d=6\ \mathrm{mm}$}、\textbf{测量段长度 $L=2\ \mathrm{m}$}（图中以 $L$ 标注测量段）；细管两端各接出一根引出管，共同接至一只 \textbf{U 形水银压差计}，压差计内盛\emph{水银}（图中以密集斜线表示），两臂水银面的\textbf{高差为 $h=30\ \mathrm{cm}$}；管内油液作层流流动，实测流量 $Q=77\ \mathrm{cm^3/s}$，油的密度 $\rho=900\ \mathrm{kg/m^3}$；不计局部损失。}
\xstarfull{4}\yiju{E4 2015 年试题计算题第 5 题（测流／测压差类，含"水银压差计读数"的处理）；E1 slide33「第五章＝考试计算题重点」；E6 教材 5.3 节哈根--泊肃叶流量公式与 U 形水银压差计换算通式 $\dfrac{p_1-p_2}{\rho g}=\left(\dfrac{\rho_{Hg}}{\rho}-1\right)h$}\nandu{难}

\item 比重为 $0.85$、动力黏度为 $0.01\ \mathrm{Pa\cdot s}$ 的润滑油，在 $d=3\ \mathrm{cm}$ 的管道中流动。每米长管道的压强降落为 $0.15\times10^{5}\ \mathrm{Pa}$，$g$ 为重力加速度。管中作层流流动，求雷诺数。\xstarfull{3}\yiju{E4 2015 年试题选择第 11 题（"雷诺数的物理意义"）；E1 slide33「第五章＝考试计算题重点」；E6 教材 5.2 节 $\Rey=\dfrac{vd}{\nu}$ 与 5.3 节 $\lambda=64/\Rey$}\nandu{中}

\item 长度 $L=1000\ \mathrm{m}$、内径 $d=200\ \mathrm{mm}$ 的普通镀锌钢管，用来输送运动黏度 $\nu=0.355\times10^{-4}\ \mathrm{m^2/s}$ 的重油，已经测得其流量 $Q=0.038\ \mathrm{m^3/s}$。求沿程损失为多少？\xstarfull{4}\yiju{E4 2022 年 818 单选第 10 题（管径减半、压强损失倍数，答 16 倍：同属"沿程损失 $\lambda$ ＋达西公式"考点）；E5 题库 \texttt{10.pdf} 选择第 10 题（沿程水头损失之比 0.886）；E1 slide33「第五章＝考试计算题重点」；E2「第五章同为主要计算考点」}\nandu{难}

\item 比重 $0.85$、$\nu=0.125\times10^{-6}\ \mathrm{m^2/s}$ 的油在粗糙度 $\Delta=0.04\ \mathrm{mm}$ 的无缝钢管中流动，管径 $d=30\ \mathrm{cm}$，流量 $Q=0.1\ \mathrm{m^3/s}$，求沿程阻力系数 $\lambda$。\xstarfull{4}\yiju{E5 题库 \texttt{10.pdf} 选择第 10 题（紊流粗糙区 $\lambda$ 只与相对粗糙度有关）；E4 2022 年 818 单选第 10 题（沿程损失随管径的变化）；E1 slide33「第五章＝考试计算题重点」；E6 教材 5.4 节尼古拉兹分区}\nandu{难}

\item 水平管路直径由 $d_1=24\ \mathrm{cm}$ 突然扩大为 $d_2=48\ \mathrm{cm}$，在突然扩大的前后各安装一测压管，读得局部阻力后的测压管比局部阻力前的测压管水柱高出 $h=1\ \mathrm{cm}$，如习题 5.11 图所示。求管中的流量 $Q$。
\tuyi{习题 5.11 图：一段\textbf{水平放置的管路}，直径由\textbf{细端 $d_1=24\ \mathrm{cm}$}（左）\textbf{突然扩大为粗端 $d_2=48\ \mathrm{cm}$}（右），水流自左向右（图中以流向箭头标注）；在突然扩大处的\textbf{前、后各装一根竖直测压管}（图中以两根并排的竖直细管表示），\textbf{扩大后的测压管液面比扩大前的测压管液面高出 $h=1\ \mathrm{cm}$}（图中以 $h$ 标注在两液面之间）。}
\xstarfull{4}\yiju{E4 2015 年试题计算题第 3 题（"水平突然缩小管路的 $d_1=15\ \mathrm{cm}$、$d_2=10\ \mathrm{cm}$……求突然缩小的水头损失"——本题 5.12 与其一字不差，本题 5.11 是与之配对的"突然扩大"同型题）；E1 slide33「第五章＝考试计算题重点」；E6 教材 5.5 节包达--卡诺公式}\nandu{难}

\item 水平突然缩小管路的 $d_1=15\ \mathrm{cm}$，$d_2=10\ \mathrm{cm}$，水的流量 $Q=2\ \mathrm{m^3/min}$。用水银测压计测得 $h=8\ \mathrm{cm}$，如习题 5.12 图所示。求突然缩小的水头损失。
\tuyi{习题 5.12 图：一段\textbf{水平放置的管路}，直径由\textbf{粗端 $d_1=15\ \mathrm{cm}$}（左）\textbf{突然缩小为细端 $d_2=10\ \mathrm{cm}$}（右），水流自左向右；在突然缩小处的前、后各接出一根引出管，共同接至一只 \textbf{U 形水银测压计}（图中以密集斜线表示工作液体水银），两臂水银面的\textbf{高差为 $h=8\ \mathrm{cm}$}。}
\xstarfull{5}\yiju{E4 2015 年试题计算题第 3 题（原文："水平突然缩小管路的 $d_1=15\ \mathrm{cm}$，$d_2=10\ \mathrm{cm}$，水的流量 $Q=2\ \mathrm{m^3/min}$。用水银测压计测得 $h=8\ \mathrm{cm}$。求突然缩小的水头损失。（8 分）"——\textbf{与本题完全相同}）；E1 slide33「第五章＝考试计算题重点」；E2「第五章同为主要计算考点，需掌握计算套路」}\nandu{中}

\item 水从直径 $d$、长 $L$ 的铅垂管路流入大气中，水箱中液面高度为 $h$，管路局部阻力可忽略，沿程阻力系数为 $\lambda$，如习题 5.13 图所示。（1）求管路起始断面 $A$ 处相对压强。（2）$h$ 等于多少时，可使 $A$ 点的压强等于大气压。
\tuyi{习题 5.13 图：一个\textbf{水箱}，箱中\textbf{液面高于箱底 $h$}（图中以竖直尺寸线自水箱内底标到液面）；箱底中心向下接出一段\textbf{铅垂管路}，\textbf{直径 $d$、长度 $L$}（图中分别以 $d$、$L$ 标注）；管路的\textbf{起始断面（与水箱内底相接处）为 $A$}，管末端为\textbf{自由出流（流入大气）}；管路局部阻力可忽略，沿程阻力系数为 $\lambda$。}
\xstarfull{5}\yiju{E4 2015 年试题计算题第 2 题（"一变直径管段 $AB$……试判断水在管中的流动方向"与本题同属"伯努利方程＋水头损失求断面压强"考点）；E1 slide16「计算考伯努利方程」；E1 slide33「第五章＝考试计算题重点」；E6 教材 5.7 节长管／短管能量方程}\nandu{难}

\item 长管输送水只计沿程损失，当 $H$、$L$ 一定，沿程损失为 $H/3$ 时管路输送功率为最大。已知 $H=127.4\ \mathrm{m}$，$L=500\ \mathrm{m}$，管路末端可用水头 $h=2H/3$，管路末端可用功率为 $1000\ \mathrm{kW}$，$\lambda=0.024$，如习题 5.14 图所示。求管路的输送流量与管路直径。
\tuyi{习题 5.14 图：一条\textbf{水平长管}（图中以水平直管表示，横向长度标为 $L=500\ \mathrm{m}$）；管道\textbf{进口断面可提供的总水头为 $H=127.4\ \mathrm{m}$}，\textbf{末端可用水头为 $h=2H/3$}（图中以竖直尺寸线分别标出 $H$ 与 $h$）；管路末端的\textbf{可用功率为 $1000\ \mathrm{kW}$}；沿程阻力系数 $\lambda=0.024$；只计沿程损失。}
\xstarfull{5}\yiju{E5 题库 \texttt{10.pdf} 计算第 2 题（"水箱中的水通过垂直管道向大气出流……沿程阻力系数 $\lambda$……求 $Q$ 与管长 $l$ 的关系"——同属"长管／短管输送功率"考点）；E1 slide33「第五章＝考试计算题重点」；E2「第五章同为主要计算考点，需掌握计算套路」；E6 教材 5.7 节简单长管第二、三类问题（$Q$ 与 $d$ 的求法）}\nandu{难}

\item 水由具有不变水位的贮水池沿直径 $d=100\ \mathrm{mm}$ 的输水管排入大气，输水管由长度 $l$ 均为 $50\ \mathrm{m}$ 的水平段 $AB$ 和倾斜段 $BC$ 组成，$h_1=2\ \mathrm{m}$，$h_2=25\ \mathrm{m}$，如习题 5.15 图所示。为了输水管在 $B$ 处的真空度不超过 $7\ \mathrm{m}$，阀门的局部阻力系数应为多少？此时管道流量 $Q$ 为多大？沿程阻力系数 $\lambda$ 取 $0.035$，不计两管相交处 $B$ 点的局部损失。
\tuyi{习题 5.15 图：左侧为一个\textbf{贮水池}，池中水位保持不变（图中以水平液面线与斜线阴影示意）；一根\textbf{输水管}水平自池中引出至\textbf{$B$ 点}（该段为\textbf{水平段 $AB$，长度 $l=50\ \mathrm{m}$}，管径 $d=100\ \mathrm{mm}$），自 $B$ 点起转为\textbf{倾斜段 $BC$（长度亦为 $l=50\ \mathrm{m}$）}向下倾斜排入大气；\textbf{$B$ 点高于水池液面 $h_1=2\ \mathrm{m}$}（图中以 $h_1$ 标注），\textbf{$C$ 端出口低于贮水池液面 $h_2=25\ \mathrm{m}$}（图中以 $h_2$ 标注）；阀门装在 $BC$ 段出口附近（靠近 $C$ 端的排出口）；沿程阻力系数 $\lambda=0.035$，不计 $B$ 点的局部损失。}
\xstarfull{5}\yiju{E4 2022 年 818 单选第 7 题（虹吸管／输水管最高处压强小于大气压，即"真空度限制"考点）；E5 题库 \texttt{10.pdf} 选择第 8 题（虹吸管最高处压强小于大气压）；E1 slide33「第五章＝考试计算题重点」；E6 教材 5.7 节短管与真空度控制}\nandu{难}

\item 如习题 5.16 图所示，要求保证自流式虹吸管中液体流量，按长管计算。试确定：（1）当 $H=2\ \mathrm{m}$，$l=44\ \mathrm{m}$，$\nu=10^{-4}\ \mathrm{m^2/s}$，$\rho=900\ \mathrm{kg/m^3}$ 时，为保证层流，$d$ 应为多少？（2）若在距进口 $l/2$ 处断面 $A$ 上的极限真空的压强水头为 $5.4\ \mathrm{m}$，输油管在上面贮油池中油面以上的最大允许超高 $z_{\max}$ 为多少？
\tuyi{习题 5.16 图：右侧为一个\textbf{贮油池}，池中油面之上为大气压；一根\textbf{自流式虹吸管}自池中伸入油面以下，向上翻过池壁后向下倾斜引出，管内液体靠\textbf{上下游液面高差 $H=2\ \mathrm{m}$} 自流（图中以 $H$ 标注在上下游液面之间）；虹吸管\textbf{总长 $l=44\ \mathrm{m}$}，管径为 $d$；在\textbf{距进口 $l/2$ 处}取断面 $A$（该断面在管路的上升段上，标高高于贮油池油面）；（2）问中给出断面 $A$ 的\textbf{极限真空压强水头为 $5.4\ \mathrm{m}$}，要求油面以上（即油面到断面 $A$ 的铅垂高差）的\textbf{最大允许超高 $z_{\max}$}；按长管计算，只计沿程损失。}
\xstarfull{4}\yiju{E4 2022 年 818 单选第 7 题（虹吸管最高处压强小于大气压＝本册"按长管计算虹吸管＋真空度限制求安装高度"考点）；E4 2015 年试题计算题第 2 题（判断流动方向）；E1 slide33「第五章＝考试计算题重点」；E6 教材 5.7 节虹吸管计算两步法}\nandu{难}

\item 两水箱之间用三根不同直径相同长度的水平管道 1、2、3 相连接，如习题 5.17 图所示。已知 $d_1=10\ \mathrm{cm}$，$d_2=20\ \mathrm{cm}$，$d_3=30\ \mathrm{cm}$，$Q_1=0.1\ \mathrm{m^3/s}$，三管沿程阻力系数相等，求 $Q_2$、$Q_3$。
\tuyi{习题 5.17 图：左侧为一个\textbf{高水位水箱}，右侧为一个\textbf{低水位水箱}，两水箱\textbf{液面高差为 $H$}；两箱之间用\textbf{三根互相平行的水平管道}连接——上管为\textbf{管道 1，$d_1=10\ \mathrm{cm}$}、中管为\textbf{管道 2，$d_2=20\ \mathrm{cm}$}、下管为\textbf{管道 3，$d_3=30\ \mathrm{cm}$}；三管\textbf{长度相同}（图中均标注为 $L$），\textbf{沿程阻力系数相等}（均为 $\lambda$）；已知管道 1 中通过的流量 $Q_1=0.1\ \mathrm{m^3/s}$；按长管计算，不计局部损失。}
\xstarfull{5}\yiju{E1 slide33「第五章＝考试计算题重点」；E2「第五章同为主要计算考点，需掌握计算套路」；E1 slide16「计算考伯努利方程」；E6 教材 5.7 节并联管路"损失相等、流量相加"与比例式 $Q_i\propto\sqrt{d_i^5/(\lambda_i l_i)}$}\nandu{难}

\item 有一直径 $d=2\ \mathrm{cm}$ 的锐缘孔口，在水头 $H=2\ \mathrm{m}$ 作用下泄流人一矩形水箱，水箱横截面积 $A=10\ \mathrm{cm}\times20\ \mathrm{cm}$，测得收缩断面直径 $d_c=1.6\ \mathrm{cm}$，经过 $49.2\ \mathrm{s}$ 后水箱内水位上升了 $3\ \mathrm{cm}$，求该孔口的收缩系数 $\varepsilon$、流量系数 $\mu$、流速系数 $\varphi$ 和局部阻力系数 $\zeta$。
\tuyi{习题 5.18 图：左侧为一个\textbf{薄壁水箱}，其侧壁下部开有一个\textbf{锐缘圆孔口，直径 $d=2\ \mathrm{cm}$}；孔口\textbf{作用水头 $H=2\ \mathrm{m}$}（图中以竖直尺寸线自水面标到孔口中心）；水由锐缘孔口射出（图中在孔口外侧画出\textbf{收缩断面，直径 $d_c=1.6\ \mathrm{cm}$}）后落入右侧的\textbf{矩形水箱内}；矩形水箱的\textbf{横截面积为 $A=10\ \mathrm{cm}\times20\ \mathrm{cm}$}；经 $49.2\ \mathrm{s}$ 后水箱内水位上升了 $3\ \mathrm{cm}$。}
\xstarfull{5}\yiju{E4 2026 回忆版计算题第 5 题（"计算管嘴体积流量（书上例题）"——同为孔口／管嘴出流类计算，教材例题与本节公式同源）；E4 2022 年 818 单选第 8 题（圆柱形直角外管嘴的正常工作条件 $H_0\le9\ \mathrm{m}$、$l=(3\sim4)d$）；E1 slide33「第五章＝考试计算题重点」；E6 教材 5.6 节薄壁小孔口出流与 $\varepsilon$、$\varphi$、$\mu$ 三系数定义}\nandu{难}

\item 在薄壁水箱上开一孔径 $d=10\ \mathrm{mm}$ 的圆孔，水箱水面位于孔口中心高度 $H=4\ \mathrm{m}$，孔口中心离地面高度 $z=5\ \mathrm{m}$，如习题 5.19 图所示。通过实验，测定射流与地面交点中心离水箱壁距离 $x=8.676\ \mathrm{m}$，孔口出流量 $Q=0.43\ \mathrm{L/s}$。如不计射流受空气的阻力，求此孔口出流的流量系数 $\mu$、流速系数 $\varphi$ 和局部阻力系数 $\zeta$。
\tuyi{习题 5.19 图：左侧为一个\textbf{薄壁水箱}，其侧壁开有一\textbf{圆孔，孔径 $d=10\ \mathrm{mm}$}；水箱内\textbf{水面位于孔口中心以上 $H=4\ \mathrm{m}$}（图中以竖直尺寸线自水箱水面标到孔口中心）；\textbf{孔口中心离地面高度 $z=5\ \mathrm{m}$}（图中以竖直尺寸线自孔口中心标到地面）；水自孔口\textbf{水平射出}，射流在重力作用下弯成抛物线，\textbf{与地面的交点中心离水箱壁的水平距离为 $x=8.676\ \mathrm{m}$}（图中以水平尺寸线标出 $x$，以竖直尺寸线分别标出 $H$ 与 $z$）；实测孔口出流量 $Q=0.43\ \mathrm{L/s}$；不计空气阻力（此即理想抛体）。}
\xstarfull{4}\yiju{E4 2026 回忆版计算题第 5 题（孔口／管嘴出流体积流量）；E4 2022 年 818 单选第 8 题（管嘴正常工作条件）；E1 slide33「第五章＝考试计算题重点」；E6 教材 5.6 节孔口出流三系数与射程关系}\nandu{难}

\item 减少水击压力的措施是什么？\xstarfull{4}\yiju{E4 2026 回忆版简答题考"水击现象"（本题即该考点的措施部分）；E1 slide33「第五章＝考试计算题重点」；E2「第五章同为主要计算考点」；E6 教材 5.8.5 节"减小水击危害的措施"}\nandu{易}

\end{ti}

\begin{tishi}
\textbf{〔本章题目范围说明〕}教材第 5 章习题自书 p95 的 5.1 起，至书 p98 的 5.20 止，共 \textbf{20 题}；
书 p99 起为第 6 章正文（已核对 \texttt{page-0108.txt} 首行为"第 6 章 相似理论与量纲分析"），
故第 5 章习题无遗漏、无误收。本章题型分布：选择题 1 组（含 11 小问）、简答／概念题 2 题（5.2、5.20）、
计算题 17 题（5.3--5.19）。
\end{tishi}

\begin{tishi}
\textbf{〔OCR 订正清单（题面部分）〕}
\begin{xiaowen}
\item \textbf{第 5.3 题：}OCR 作"流通运动黏度 $\nu=1.306\times10^{-6}\ \mathrm{m^2/s}$"，"流通"应为"水的"（原题指管内水流的运动黏度），已订正为 $\nu=1.306\times10^{-6}\ \mathrm{m^2/s}$。
\item \textbf{第 5.4 题：}OCR 把题目编号排成"$(R^2-r^2)$，其中……"并与 5.3 混行，且把 $u$ 印成"流速 $v=$"；按教材 5.3 节口径订正为 $u=\dfrac{\Delta p}{4\mu L}\left(R^2-r^2\right)$（$u$ 为点流速，$v$ 为断面平均流速，两者不可混用）。
\item \textbf{第 5.5 题：}OCR 作"流量 $Q=1.0\ \mathrm{cm^2/s}$"，量纲错误，按流量单位订正为 $Q=1.0\ \mathrm{cm^3/s}=1.0\times10^{-6}\ \mathrm{m^3/s}$。
\item \textbf{第 5.6 题：}OCR 作"$\rho=0.85\ \mathrm{g/cm^2}$"，量纲错误，密度单位应为 $\mathrm{g/cm^3}$，订正为 $\rho=0.85\ \mathrm{g/cm^3}=850\ \mathrm{kg/m^3}$；OCR 作"$\nu=0.18\ \mathrm{cm^2/s}$"，尺寸为长度平方除以时间，量纲正确（$1\ \mathrm{cm^2/s}=10^{-4}\ \mathrm{m^2/s}$），保留原值；OCR 的"管壁切应力 $t_0$"订正为 $\tau_0$。
\item \textbf{第 5.7 题：}OCR 作"流量 $Q=77\ \mathrm{cm^2/s}$"，量纲错误，按流量单位订正为 $Q=77\ \mathrm{cm^3/s}=77\times10^{-6}\ \mathrm{m^3/s}$；OCR 作"油密度 $\sigma=900\ \mathrm{kg/m^3}$"（且末尾多一个引号），按密度符号订正为 $\rho=900\ \mathrm{kg/m^3}$。
\item \textbf{第 5.8 题：}OCR 作"动力黏度为 $0.01\ \mathrm{gPa\cdot s}$"，符号不清；按动力黏度的量级与润滑油的通常取值，原文应为 $\mu=0.01\ \mathrm{Pa\cdot s}$，已按此处理（$\mathrm{Pa\cdot s}$ 是 SI 单位，加"g"则量级荒谬）。
\item \textbf{第 5.9 题：}OCR 作"运动黏度 $\nu=0.355\times10^{-}\ \mathrm{m^2/s}$"，指数缺失；按重油运动黏度 $\nu=0.355\times10^{-4}\ \mathrm{m^2/s}$（$=3.55\times10^{-5}\ \mathrm{m^2/s}$，约 $35.5\ \mathrm{mm^2/s}$，与重油量级相符）处理；流量 OCR 作 $q=0.038\ \mathrm{m^2/s}$，量纲错误，订正为 $Q=0.038\ \mathrm{m^3/s}$。
\item \textbf{第 5.10 题：}OCR 作"$\nu=0.125\times10^{-}\ \mathrm{m^2/s}$"，指数缺失。按\textbf{参考答案反推}：教材答案为 $\lambda\approx0.038$，由布拉修斯式 $\lambda=0.3164/\Rey^{0.25}$ 解得 $\Rey=(0.3164/0.038)^{4}=4.8\times10^3$，对应 $\nu=vd/\Rey=1.4147\times0.3/4800=8.8\times10^{-5}\ \mathrm{m^2/s}$，即\textbf{原书应为 $\nu=0.125\times10^{-4}\ \mathrm{m^2/s}$}（此时 $\Rey=3395$ 属湍流光滑区，$\lambda=0.0414$，与参考答案 $0.038$ 同量级）。若照 OCR 字形取 $0.125\times10^{-6}$，则 $\Rey=3.40\times10^6$ 进入湍流过渡区、$\lambda\approx0.013$，与参考答案相差约 3 倍。本册在题面\textbf{保留 OCR 原字形}，解答中把两种口径的完整结果都算出，读者以自己书上印的指数为准。另：OCR 作"粗糙度 $=0.04\ \mathrm{mm}$"，补为绝对粗糙度 $\Delta=0.04\ \mathrm{mm}$；OCR 两次出现的 "$q=0.1\ \mathrm{m^2/s}$"（含 5.17 题）量纲错误，流量单位应为 $\mathrm{m^3/s}$。
\item \textbf{第 5.14 题：}OCR 作"沿程损失为 $H/3$ 时管路输送功率为最大"，并给出"末端可用水头 $h=2H/3$"；两处均可读，按原文收录。
\item \textbf{第 5.15 题：}OCR 作"真空度不超过 $7\ \mathrm{m^2}$"，量纲错误，真空度是水头（长度），订正为 $7\ \mathrm{m}$（水柱）。
\item \textbf{第 5.16 题：}OCR 作"$\nu=10^{-}$"，指数缺失；取 $\nu=10^{-4}\ \mathrm{m^2/s}$（$=100\ \mathrm{mm^2/s}$，与题给的油品量级相符），已按此处理；OCR 的"$\sigma=900\ \mathrm{kg/m^2}$"按密度符号与单位订正为 $\rho=900\ \mathrm{kg/m^3}$；OCR 的"最大允许超高 $\max$"补为 $z_{\max}$。
\item \textbf{第 5.18 题：}OCR 作"水箱横截面积 $A=10\ \mathrm{cm}\times20\ \mathrm{cm}$"，但该量是\textbf{面积}、$A$ 与第五章的过流面积符号重复，本册在题面保留原符号而在解答中写清"矩形水箱横截面积"以免与孔口面积混淆；OCR 的"收缩断面直径 $d=1.6\ \mathrm{cm}$"补为 $d_c=1.6\ \mathrm{cm}$（与原孔径 $d=2\ \mathrm{cm}$ 相区别）；OCR 的"收缩系数 $\in$"订正为 $\varepsilon$，并补出漏印的"流速系数 $\varphi$"（原题四个待求量按教材 5.6 节为 $\varepsilon$、$\mu$、$\varphi$、$\zeta$）。
\item \textbf{第 5.19 题：}OCR 作"孔口出流量 $Q=0.43\ \mathrm{L/s}$"（$=0.43\times10^{-3}\ \mathrm{m^3/s}$），已按流量单位处理；OCR 的"流速系数"后缺符号，按教材 5.6 节补为 $\varphi$。
\item \textbf{第 5.1 题第（11）小问：}OCR 仅排出一个"A."，其余选项缺字母；按教材原题补全 A--D 四个选项。
\item \textbf{第 5.1 题第（8）（9）小问：}OCR 的"$L_2=3L$"补为 $L_2=3L_1$；"$Q=Q_2$""$Q=1.5Q_2$"等缺下标，按题意补为 $Q_1$、$Q_2$；"$Q=0$"按题意补为 $Q_2=0$。
\item \textbf{第 5.2 题：}OCR 原文"两种液体阻力及能量损失形式和它们的计算公式分别是什么？"——按教材 5.1 节口径，"液体阻力"应为"流体阻力"（本册在题面保留 OCR 字形，解答中按"流体阻力"作答并说明）。
\end{xiaowen}
\end{tishi}

\begin{tishi}
\textbf{〔存疑处（做题时请以题图为准）〕}
\textbf{第 5.13 题：}$A$ 点定在"管路起始断面（与水箱内底相接处）"由题图与题意推出；
若你的题图中 $A$ 点画在箱底以下的某一距离处，请把该距离并入 $L$，方法与解答完全相同。\\
\textbf{第 5.15 题：}阀门的安装位置（$BC$ 段靠近出口）由题图的阀门符号读出；
$h_1=2\ \mathrm{m}$、$h_2=25\ \mathrm{m}$ 的基准分别为"$B$ 点到水池液面"与"$C$ 端到水池液面"，
若你的题图尺寸线基准不同，请按图中实际尺寸线取值。\\
\textbf{第 5.16 题：}第（1）问"为保证层流，$d$ 应为多少"按"$d$ 不超过层流允许的最大直径"理解（取 $\Rey=2000$ 为界）；
第（2）问的"极限真空压强水头 $5.4\ \mathrm{m}$"是断面 $A$ 处的允许最大真空值。两问所用 $\nu$ 的指数由 OCR 缺失处补出，
若你的书上印的是其他指数，请按书中值代入，方法与解答不变。\\
\textbf{第 5.18 题：}"经过 $49.2\ \mathrm{s}$ 后水箱内水位上升了 $3\ \mathrm{cm}$"是测流量的实验记录，
解答把它换算为孔口的平均出流量（水箱很浅、$49.2\ \mathrm{s}$ 内水位只升 $3\ \mathrm{cm}$，按定水头处理误差很小）。
\end{tishi}
